\documentclass[11pt]{article}
\usepackage[margin=1in]{geometry}
\usepackage{amsmath,amssymb,booktabs,hyperref}
\newcommand{\NS}{\mathrm{NS}}
\newcommand{\eps}{\boldsymbol\epsilon}
\newcommand{\bet}{\boldsymbol\beta}
\newcommand{\hp}{\boldsymbol h}
\newcommand{\pp}{\boldsymbol p}
\newcommand{\Pol}{\operatorname{Pol}}
\title{External NS upper components in the sphere pillow $h$-recursion}
\author{Machine note: human-note conventions only}
\date{30 August 2026}
\begin{document}
\maketitle

\paragraph{Status and scope.}
This extends the bottom-component research code to external
$G_{-1/2}\phi$ states of intrinsically even NS superprimaries.  We give
the component conformal factor, parity transport, and fixed-$c$ pole
recursion.  For general cap components the regular term is a polynomial
in the common internal weight, not the bottom-component unit seed.
An exact, finite-cutoff PBW polynomial-division prescription computes
these regular terms and gives an executable reference recursion.  It is
not a new fast all-orders formula for the general seed.  A simpler
unit-seed extension for upper components confined to mobile insertions
has a separate power-counting argument and independent numerical tests.
The human note and the released bottom-component package are unchanged.
No production amplitude or subtraction prescription is changed here.

\section{States and three-point functions}

Write
\begin{equation}
 V_j^{(\beta_j)}=(G_{-1/2})^{\beta_j}\phi_{d_j},\qquad
 \beta_j\in\{0,1\},\qquad D_j=d_j+\frac{\beta_j}{2}.
\end{equation}
The $d_j$ remain the superprimary weights.  The $D_j$ are the weights
under ordinary conformal transformations: $G_{-1/2}\phi_d$ is a
Virasoro primary because $[L_m,G_{-1/2}]\phi_d=0$ for $m>0$.
It is not a new NS superprimary.

We use exactly
\begin{equation}
 \rho_0(\phi_1,\phi_2,\phi_3)=1,\qquad
 \rho_1(\phi_1,G_{-1/2}\phi_2,\phi_3)=1.
\end{equation}
For even intrinsic primary parities the human-note table is
\begin{center}
\begin{tabular}{ccl}
\toprule
$(a_1a_2a_3)$ & $a=\sum a_i\bmod2$ &
$\rho_a(G_{-1/2}^{a_1}\phi_{h_1},G_{-1/2}^{a_2}\phi_{h_2},G_{-1/2}^{a_3}\phi_{h_3})$\\
\midrule
000&0&$1$\\
110&0&$h_1+h_2-h_3$\\
101&0&$h_1-h_2+h_3$\\
011&0&$h_1-h_2-h_3$\\
100&1&$1$\\
010&1&$1$\\
001&1&$-1$\\
111&1&$-(h_1+h_2+h_3-\tfrac12)$\\
\bottomrule
\end{tabular}
\end{center}
Thus the upper components have their own three-point numerators; simply
feeding $D_j$ into the bottom-component fusion polynomials is incorrect.
Nor should these three-point coefficients be divided out to impose an
artificial unit level-zero normalization.

The Ward engine uses, for $W_d=G_{-1/2}\phi_d$,
\begin{align}
 [G_r,\phi_d(z)]&=z^{r+1/2}W_d(z),\\
 \{G_r,W_d(z)\}&=z^{r-1/2}
       \bigl(z\partial_z+(2r+1)d\bigr)\phi_d(z),\\
 [L_m,V^{(\beta)}_d(z)]&=z^m
       \bigl(z\partial_z+(m+1)(d+\beta/2)\bigr)V^{(\beta)}_d(z).
\end{align}
The signs are those of the fixed-parity three-form, not an extra BRY
scalar rephasing.  External upper states at zero and infinity are
explicit ket and bra $G_{-1/2}$ states in the direct PBW implementation.

\section{The sphere--pillow factor and component labels}

Let $m=n-3$ and place the ordered fields at
$(0,z,t_1,\ldots,t_{m-1},1,\infty)$.  Keep
\begin{equation}
 c=\frac32+3(b+b^{-1})^2,\qquad
 q=e^{i\pi\tau}=\prod_{i=1}^m p_i,\qquad
 \tau=i\frac{K(1-z)}{K(z)},\quad
 K(z)={}_2F_1(\tfrac12,\tfrac12;1;z).
\end{equation}
The Abel--Jacobi map and the collision-selected segment nomes are
unchanged.  In particular
\begin{equation}
 (\varrho_1,\ldots,\varrho_m)=
 \begin{cases}(16q),&m=1,\\
 (4p_1,p_2,\ldots,p_{m-1},4p_m),&m\ge2.
 \end{cases}
\end{equation}
All powers below initially use the ordered real OPE sheet.  A complex
continuation must carry consistent spin and logarithm lifts.

Use the same normalization of the elliptic block as in the bottom note:
\begin{equation}\label{eq:reconstruction}
 \boxed{\displaystyle
 \mathcal F_{\bet,\eps}=
 \Lambda_n^{(c)}(D_1,\ldots,D_n)
 \prod_{i=1}^m\varrho_i^{h_i-c/24}\,
 C_{\NS}(q)\,H_{\bet,\eps}(\hp;\pp).}
\end{equation}
Here
\begin{align}
 \Lambda_n^{(c)}(D)
 &=z^{c/24-D_1-D_2}(1-z)^{c/24-D_2-D_{n-1}}
   \theta_3(q)^{c/2-4(D_1+D_2+D_{n-1}+D_n)-2\sum_{j=3}^{n-2}D_j}
   \notag\\
 &\quad\times\prod_{j=1}^{m-1}[t_j(1-t_j)(t_j-z)]^{-D_{j+2}/2},\\
 C_{\NS}(q)&=\theta_3(q^2)\prod_{k\ge1}(1-q^{2k})^{-3/4}.
\end{align}
The actual central charge is used, and $C_{\NS}$ is held fixed as a
normalization factor.  For cap upper components it is not by itself the
whole regular part.

Let $\epsilon_i$ be the parity of the internal level on edge $i$.
The relative three-form labels are now
\begin{align}\label{eq:labels}
 \alpha_1&=\epsilon_1+\beta_1+\beta_2,\\
 \alpha_v&=\epsilon_{v-1}+\epsilon_v+\beta_{v+1},\qquad 2\le v\le m,\\
 \alpha_{m+1}&=\epsilon_m+\beta_{n-1}+\beta_n
 \qquad\pmod 2.
\end{align}
In particular $\sum_v\alpha_v=\sum_j\beta_j\pmod2$.  Every external
marking still has $2^m$ internal parity sectors.  They are not to be
summed without the corresponding CFT structure constants.

\section{Pole term and common-weight recursion}

Define
\begin{equation}
 h_{rs}=\frac{(b+b^{-1})^2-(rb+s/b)^2}{8},\qquad
 \ell_{rs}=\frac{rs}{2},\qquad \pi_{rs}=rs\bmod2,
 \quad r,s\ge1,\quad r+s\in2\mathbb Z.
\end{equation}
The $(1,1)$ pole is included.  For a pole on edge $k$, put
\begin{equation}
 \widehat h_j=h_j-h_k+h_{rs},\qquad
 h'_j=\widehat h_j+\ell_{rs}\delta_{jk},\qquad
 \epsilon'_j=\epsilon_j+\pi_{rs}\delta_{jk}\pmod2.
\end{equation}
The external markings do not change.  Exactly the two incident
three-form labels toggle for an odd null vector.

With the human-note inverse norm slope $A_{rs}^c$ and fusion
polynomials $P_{rs}^{\alpha}(x,y;c)$, the residue kernel is
\begin{equation}\label{eq:R}
 R_{rs}^{(k;\bet,\eps)}
 =(-1)^{rs}A_{rs}^c
 P_{rs}^{\alpha_k}(L_k;c)P_{rs}^{\alpha_{k+1}}(R_k;c),
\end{equation}
where the ordered weight pairs are
\begin{align}
 L_k&=\begin{cases}(d_1,d_2),&k=1,\\
              (\widehat h_{k-1},d_{k+1}),&k>1,
       \end{cases}\\
 R_k&=\begin{cases}(d_n,d_{n-1}),&k=m,\\
              (\widehat h_{k+1},d_{k+2}),&k<m.
       \end{cases}
\end{align}
Crucially, these use $d_j$, not $D_j$.  The component dependence is
carried by the labels, the unchanged upper-state insertions in the child
block, and the regular seed.  Null-vector factorization is an identity
for arbitrary descendant states in the other slots, so it applies also
when those states are external $G_{-1/2}\phi$.

For completeness, with $\lambda_x^2=(b+b^{-1})^2-8x$,
\begin{align}
 A_{rs}^c&=\frac12
  \prod_{\substack{u=1-r,\ldots,r;\ v=1-s,\ldots,s\\u+v\ {
m even}\\
                    (u,v)\ne(0,0),(r,s)}}
           \left(\frac{ub+v/b}{\sqrt2}\right)^{-1},\\
 P_{rs}^{\alpha}(x,y;c)&=
  \prod_{\substack{u=1-r,3-r,\ldots,r-1\\v=1-s,3-s,\ldots,s-1\\
             u+v-r-s=2(1-\alpha)\ ({\rm mod}\ 4)}}
   \frac{\lambda_x-\lambda_y+ub+v/b}{2\sqrt2}
   \frac{\lambda_x+\lambda_y+ub+v/b}{2\sqrt2}.
\end{align}
No new fusion-polynomial normalization is introduced.

Write $h_1=H$, $h_j=H+a_j$, with $a_1=0$.  The resulting recursion is
\begin{equation}\label{eq:recursion}
 \boxed{\displaystyle
 H_{\bet,\eps}(\hp;\pp)
 =S_{\bet,\eps}(H,\boldsymbol a;\pp)
 +\sum_{k=1}^m\sum_{\substack{r,s\ge1\\r+s\ {
m even}}}
   \frac{\varrho_k^{rs/2} R_{rs}^{(k;\bet,\eps)}}{h_k-h_{rs}}
   H_{\bet,\eps'}(\hp';\pp).}
\end{equation}
Differences are recomputed in each child.  They are not frozen to their
initial values after shifting only the null edge.

\section{The regular part: why a unit seed is not general}

At a fixed coefficient $\pp^{\boldsymbol N/2}$ the block is a rational
function of the internal weights.  Its coordinate pullback and division
by the prefactor in \eqref{eq:reconstruction} add no internal-weight
poles.  Therefore, away from confluent parameters, subtracting the
simple principal parts leaves a polynomial in $H$ at fixed differences.
Define the seed coefficient unambiguously by
\begin{equation}\label{eq:seed}
 S_{\bet,\boldsymbol N}(H,\boldsymbol a)=
 \Pol_{H=\infty}
 \left[H_{\bet,\boldsymbol N}(H,H+a_2,\ldots,H+a_m)\right],
 \qquad
 S_{\bet,\eps}=\sum_{\boldsymbol N\bmod2=\eps}
       S_{\bet,\boldsymbol N}\,\pp^{\boldsymbol N/2}.
\end{equation}
$\Pol$ retains \emph{all} nonnegative Laurent powers, including the
constant.  It does not mean retaining only the highest power of $H$.

An exact finite-order computation of \eqref{eq:seed} is:
\begin{enumerate}
 \item Sew the component three-forms with the NS PBW inverse Gram matrices.
 \item Pull the plane ratios $(z/t_1,t_1/t_2,\ldots,t_{m-1})$ back to the
 elliptic $p_i$, using $D_j$ in the conformal factor.
 \item Substitute $h_i=H+a_i$ and cancel each rational coefficient.
 \item Divide its numerator polynomial by its denominator polynomial in
 $H$; retain the quotient.
\end{enumerate}
This determines every regular coefficient at a prescribed finite order
without an asymptotic fit or guessed seed.  These polynomials can be
compiled once for fixed $c,d_j,\bet$ and reused at varying internal
weights.  Using PBW to construct them, however, means the resulting
recursion comparison is a pole-factorization test, not independent
evidence for a proposed closed seed.  It is not yet a production speedup
for arbitrary cap components.

\subsection{Four-point counterexample to the unchanged seed}

For $\bet=(0,1,1,0)$, the odd leading plane coefficient is
\begin{equation}
 B_{1/2}=-\frac{(h+d_2-d_1)(h+d_3-d_4)}{2h}.
\end{equation}
Because $(16q)^{1/2}=4q^{1/2}$, its elliptic regular part starts as
\begin{align}
 S_{0110,0}&=1+O(q^3),\\
 S_{0110,1}&=
   2(-h+d_1-d_2-d_3+d_4)q^{1/2}
   +4(2d_1-2d_2-1)(2d_3-2d_4+1)q^{3/2}
   +O(q^{5/2}).
\end{align}
The displayed integer seed coefficients at $q$ and $q^2$ vanish.
The linear term in $h$ is required by the same human-note three-form
normalization, not an optional change of convention.

\subsection{Five points with three upper components}

For $\bet=(0,1,1,1,0)$ and $a=h_2-h_1$, the regular coefficients through
total degree one are
\begin{align}
 S_{\bet,00}&=1+O_{>1},\\
 S_{\bet,10}&=(a-d_3)p_1^{1/2}+O_{>1},\\
 S_{\bet,01}&=-(a+d_3)p_2^{1/2}+O_{>1},\\
 S_{\bet,11}&=
 \left[2h_1+a+d_3-\frac12+2(d_2+d_4-d_1-d_5)\right]
       (p_1p_2)^{1/2}+O_{>1}.
\end{align}
Here $O_{>1}$ denotes total physical degree greater than one.
The integer degree-one terms in the $00$ sector vanish.
All four parity sectors must be kept as separate blocks.

\section{Fast special case: upper components only at mobile insertions}

Suppose
\begin{equation}\label{eq:bottom-caps}
 \beta_1=\beta_2=\beta_{n-1}=\beta_n=0,
\end{equation}
with arbitrary markings at $t_j$.  There is a useful extension of the
previous large-common-weight argument.  At fixed PBW levels and fixed
$h_{i+1}-h_i$, the commutators of a mode with a light insertion bring
only finite differences and light weights.  Commuting matched internal
modes through the insertion gives the leading Gram contraction.  For
$W=G_{-1/2}\phi$, moving a fermionic mode through $W$ contributes a
minus sign, so the normalized leading diagonal matrix element differs
from that for $\phi$ by $(-1)^F$.  Off-diagonal transitions are suppressed
in the same oscillator power counting.

The bottom caps in the original proposal retain only the all-even
internal parity sector at leading order.  Hence $F$ is even on every
surviving propagator and the additional diagonal signs are all $+1$.
Under the same coefficientwise cap-asymptotic assumptions as the
bottom-component proposal, this gives
\begin{equation}\label{eq:interior-seed}
 S_{\bet,\eps}=\delta_{\eps,0}\qquad\text{when \eqref{eq:bottom-caps} holds}.
\end{equation}
Before dividing by $C_{\NS}$, the regular term is
$\delta_{\eps,0}C_{\NS}(q)$.  This is a proposed all-orders conclusion
of the shared power-counting assumptions; the finite numerical tests
below are not a proof of uniform convergence in the moduli.

For five points this covers $(0,0,1,0,0)$, namely a single upper
component at the mobile insertion.  The corresponding $h$ evaluator
uses no PBW or $c$-recursion coefficients.  This fast seed is explicitly
rejected by the code when any of the four cap fields is upper.

\section{Implementation and validation boundaries}

The research implementation is
\begin{quote}\ttfamily\small
Code/sphere\_five\_kummer\_h\_recursion/\\
ns\_pillow\_components.py
\end{quote}
It supplies \texttt{ComponentEllipticRecursion},
\texttt{ExactPBWSeeds}, compiled numerical polynomial seeds, and the
explicitly selected \texttt{InteriorUnitSeed}.
The direct PBW engine now accepts the option
\texttt{external\_\allowbreak descendants}; its all-bottom behavior is
unchanged.  Missing general upper-component seeds raise an error.

The accompanying \texttt{check\_ns\_pillow\_components.py} separates:
the eight exact human-note three-forms; independent component PBW versus
$c$-recursion; symbolic pole-factorization certificates with PBW-derived
regular polynomials; and independent tests of \eqref{eq:interior-seed}.
Detailed results, parameter scope, and reproducible commands are recorded
in the companion Markdown note and JSON ledgers.

At total degree three, all 16 four-point and all 32 five-point component
markings pass the pole certificates with arbitrary internal weights and
fixed exact external parameters.  In addition, the four-point marking
$0110$ and five-point marking $01110$ pass with $b$ and every external
and internal weight symbolic; the latter includes 28 coefficients and
106 residue identities.  Independently of these PBW-derived seeds,
the mobile-upper five-point marking $00100$ agrees with PBW and
$c$-recursion through total degree ten for four generic numerical
fixtures.  These are coefficient checks, not amplitude-level tests.

The present extension does not implement a confluent $b=1$ limit,
complex-moduli spin transport, or the physical PCO component sum.
General cap-component seed polynomials are available at a chosen finite
cutoff, but a compact all-orders seed formula remains a separate task.
Only after those pieces and amplitude-level comparisons are complete
should the proposed $|q|<0.3$ hybrid backend be enabled.  Polynomial
subtraction remains above that eventual block-evaluation layer.

\end{document}
